\documentclass[twocolumn]{article}
\usepackage[margin=1in]{geometry}
\usepackage[utf8]{inputenc}
\usepackage[T1]{fontenc}
\usepackage{hyperref}
\usepackage{url}
\usepackage{booktabs}
\usepackage{multirow}
\usepackage{amsfonts}
\usepackage{amsmath}
\usepackage{amssymb}
\usepackage{amsthm}
\usepackage{graphicx}
\usepackage{xcolor}
\usepackage{subcaption}
\usepackage{natbib}
\usepackage{microtype}

%% 9pt body text to approximate PNAS style
\usepackage{setspace}
\renewcommand{\normalsize}{\fontsize{9}{11}\selectfont}
\normalsize

%% Theorem environments
\newtheorem{theorem}{Theorem}
\newtheorem{lemma}[theorem]{Lemma}
\newtheorem{corollary}[theorem]{Corollary}
\newtheorem{definition}[theorem]{Definition}
\newtheorem{proposition}[theorem]{Proposition}
\newtheorem{remark}[theorem]{Remark}

%% Macros
\newcommand{\R}{\mathbb{R}}
\newcommand{\E}{\mathbb{E}}
\newcommand{\Var}{\mathrm{Var}}
\newcommand{\SHAP}{\textsc{SHAP}}
\newcommand{\DASH}{\textsc{Dash}}
\newcommand{\GBDT}{\textsc{GBDT}}
\newcommand{\Params}{\mathbf{\Theta}}
\newcommand{\Hyp}{\mathcal{H}}
\newcommand{\Obs}{\mathcal{O}}
\newcommand{\ExplSpace}{\mathcal{E}}
\newcommand{\sE}{\mathcal{E}}
\newcommand{\sH}{\mathcal{H}}
\newcommand{\sO}{\mathcal{O}}
\newcommand{\incomp}{\perp}

\title{The Impossibility of Faithful Explanation Under Underdetermination}

\author{
  Drake Caraker\thanks{Corresponding author: \texttt{drakecaraker@gmail.com}} \quad
  Bryan Arnold \quad
  David Rhoads \\[4pt]
  \textit{Independent Researchers}
}

\date{}

\begin{document}

\twocolumn[
  \begin{@twocolumnfalse}
  \maketitle

%%=======================================================================
\begin{abstract}
%%=======================================================================
When a system is underspecified---when multiple configurations produce
the same observable output---any explanation faces a fundamental
tradeoff. We prove that no explanation can simultaneously be faithful
(consistent with the system's native explanation), stable (invariant
across observationally equivalent configurations), and decisive
(committing to every distinction the native explanation makes). This
impossibility holds whenever the Rashomon property obtains: equivalent
configurations with incompatible explanations. We prove the Rashomon
property is both necessary and sufficient---the exact boundary between
possibility and impossibility. The theorem is instantiated across nine
domains: feature attribution, attention maps, counterfactual
explanations, concept probes, causal discovery, model selection,
saliency maps, language model explanations, and mechanistic
interpretability. For causal discovery, the Rashomon property is derived
from first principles (Markov equivalence). The resolution---$G$-invariant
aggregation---is proved uniquely optimal. The result is mechanically
verified in the Lean~4 proof assistant (346 theorems, 0 unproved goals).
\end{abstract}

\medskip

\noindent\textbf{Significance.}\quad
Explainability is central to trustworthy AI, scientific inference, and
regulatory compliance. We prove a universal impossibility: no method of
explaining an underspecified system can simultaneously be faithful to the
system, stable across equivalent configurations, and decisive in its
commitments. This formalizes the century-old Quine--Duhem underdetermination
thesis as a mathematical theorem for the first time. The impossibility
applies to every major explanation paradigm in machine learning---from
SHAP values to neural network circuit analysis---and extends to scientific
explanation, medical diagnosis, and legal interpretation. The constructive
resolution (aggregate over equivalent configurations) provides the
mathematically optimal tradeoff. The entire proof is mechanically verified,
eliminating logical gaps.

\medskip

\noindent\textbf{Keywords:} impossibility theorem, explainability, underdetermination, Rashomon effect, formal verification

\vspace{1em}
\end{@twocolumnfalse}
]

%%=======================================================================
\section{Introduction}
\label{sec:intro}
%%=======================================================================

When multiple configurations of a system produce identical observable
output, any attempt to explain the system's behavior faces a fundamental
constraint. This paper proves that constraint is an impossibility: no
explanation can simultaneously be faithful, stable, and decisive.

This structure arises across science. In philosophy of science,
Quine~\cite{quine1951dogmas} and Duhem~\cite{duhem1906aim} argued that
empirical evidence underdetermines theory choice---no finite body of
observations can uniquely select a scientific theory from among
empirically adequate alternatives. In neuroscience, Marder and
Goaillard~\cite{marder2006variability} documented that nervous systems
with identical functional output are built from qualitatively different
circuit configurations, a phenomenon called \emph{neural degeneracy}
\cite{edelman2001degeneracy}. In statistics and machine learning,
Breiman's Rashomon effect~\cite{breiman2001random} and its modern
elaboration by Rudin et al.~\cite{rudin2024amazing} describe the
multiplicity of near-optimal models with divergent internal structure. In
economics, Manski~\cite{manski2003partial} developed partial
identification theory precisely because observational data generically
fail to pin down a unique model. In database theory,
Bancilhon and Spyratos~\cite{bancilhon1981update} proved that view updates
are inherently ambiguous when the view is not injective.

We prove the impossibility, characterize its exact boundary (the Rashomon
property---necessary and sufficient), demonstrate it across nine domains
with the most detailed treatment in machine learning, and provide a
constructive resolution. Informally:

\begin{quote}
\emph{If two configurations of a system produce the same observable output
but honestly report incompatible explanations, then no explanation method
can simultaneously be faithful to the system, stable across equivalent
configurations, and decisive in its commitments.}
\end{quote}

\noindent
The proof is four steps---decisiveness, stability, faithfulness,
contradiction---and is mechanically verified in the Lean~4 proof
assistant~\cite{demoura2021lean4} (73 files, 346 theorems, 72 axioms,
0~\texttt{sorry}). The core impossibility requires zero axioms beyond
the Rashomon property itself.

Our contributions are: (1)~a universal impossibility theorem unifying
nine domain-specific results under a single proof; (2)~a proof that the
Rashomon property is both necessary and sufficient---the exact boundary
between possibility and impossibility; (3)~nine instances spanning
feature attribution, attention maps, counterfactual explanations, concept
probes, causal discovery (with Rashomon \emph{derived} from Markov
equivalence), model selection, saliency maps, language model explanations,
and mechanistic interpretability; (4)~a constructive resolution via
$G$-invariant aggregation proved uniquely optimal; (5)~empirical
validation across seven instances; and (6)~complete mechanical
verification. The result connects to Hunt--Stein invariant decision
theory~\cite{lehmann2005testing} and, to our knowledge, provides the
first mathematical formalization of the Quine--Duhem underdetermination
thesis.


%%=======================================================================
\section{Framework}
\label{sec:framework}
%%=======================================================================

\subsection{Explanation Systems}

We introduce a minimal abstract framework in which the impossibility can
be stated and proved. All nine instances in Section~\ref{sec:instances}
are special cases.

\begin{definition}[Explanation System]
\label{def:explanation-system}
An \emph{explanation system} is a tuple
$\mathcal{S} = (\Params, \sH, Y, \mathsf{observe}, \mathsf{explain}, \incomp)$
where:
\begin{itemize}
  \item $\Params$ is a \emph{configuration space}---the set of possible
    configurations of the system (e.g., model parameters, experimental
    setups, circuit wirings);
  \item $\sH$ is an \emph{explanation space}---the set of possible
    explanations;
  \item $Y$ is an \emph{observable space}---the set of measurable
    outputs;
  \item $\mathsf{observe}: \Params \to Y$ assigns to each configuration
    its observable behavior;
  \item $\mathsf{explain}: \Params \to \sH$ assigns to each
    configuration the explanation it ``honestly reports''; and
  \item $\incomp \subseteq \sH \times \sH$ is an \emph{incompatibility
    relation} satisfying \emph{irreflexivity}:
    $\neg\,\incomp(h,h)$ for every $h \in \sH$.
\end{itemize}
\end{definition}

An \emph{explanation map} $E : \Params \to \sH$ is a procedure that a
practitioner \emph{chooses}. The native explanation
$\mathsf{explain}(\theta)$ records what the configuration actually
implies; $E(\theta)$ is what the chosen method reports.

\subsection{Three Desiderata}

\begin{definition}[Faithful, Stable, Decisive]
\label{def:desiderata}
An explanation map $E$ is:
\begin{itemize}
  \item \textbf{Faithful} if $E(\theta)$ never contradicts the native
    explanation:
    $\forall\,\theta,\;\neg\,\incomp(E(\theta), \mathsf{explain}(\theta))$.
  \item \textbf{Stable} if equivalent configurations receive the same
    explanation:
    $\forall\,\theta_1, \theta_2,\;
    \mathsf{observe}(\theta_1) = \mathsf{observe}(\theta_2)
    \Rightarrow E(\theta_1) = E(\theta_2)$.
  \item \textbf{Decisive} if $E(\theta)$ inherits every incompatibility
    of the native explanation:
    $\forall\,\theta, h,\;
    \incomp(\mathsf{explain}(\theta), h)
    \Rightarrow \incomp(E(\theta), h)$.
\end{itemize}
\end{definition}

\noindent
Faithfulness says the method does not contradict the system. Stability says
the method is reproducible---it does not depend on which equivalent
configuration happens to be examined. Decisiveness says the method is at
least as committal as the system itself: whatever the native explanation
rules out, the chosen explanation also rules out.

\subsection{The Rashomon Property}

\begin{definition}[Rashomon Property]
\label{def:rashomon}
An explanation system has the \emph{Rashomon property} if there exist
$\theta_1, \theta_2 \in \Params$ with
\[
  \mathsf{observe}(\theta_1) = \mathsf{observe}(\theta_2)
  \quad\text{and}\quad
  \incomp(\mathsf{explain}(\theta_1), \mathsf{explain}(\theta_2)).
\]
\end{definition}

\noindent
Two configurations produce the same observable output yet honestly report
incompatible explanations. This is the formal statement of
underdetermination at the explanation level.

\subsection{Main Theorem}

\begin{theorem}[Universal Explanation Impossibility]
\label{thm:universal}
If an explanation system has the Rashomon property, then no explanation
map $E$ can simultaneously be faithful, stable, and decisive.
\end{theorem}

\begin{proof}
Let $\theta_1, \theta_2$ witness the Rashomon property.
Suppose $E$ is faithful, stable, and decisive.
\begin{enumerate}
  \item \textbf{Decisiveness} at $\theta_1$:
    $\incomp(\mathsf{explain}(\theta_1), \mathsf{explain}(\theta_2))$
    implies $\incomp(E(\theta_1), \mathsf{explain}(\theta_2))$.
  \item \textbf{Stability}:
    $\mathsf{observe}(\theta_1) = \mathsf{observe}(\theta_2)$
    implies $E(\theta_1) = E(\theta_2)$, so
    $\incomp(E(\theta_2), \mathsf{explain}(\theta_2))$.
  \item \textbf{Faithfulness} at $\theta_2$:
    $\neg\,\incomp(E(\theta_2), \mathsf{explain}(\theta_2))$.
  \item Steps 2 and 3 contradict. \qedhere
\end{enumerate}
\end{proof}

\begin{proposition}[Tightness]
\label{prop:tightness}
Each pair of properties is simultaneously achievable:
(i)~faithful $+$ decisive (set $E = \mathsf{explain}$; stability fails);
(ii)~faithful $+$ stable (set $E$ to a neutral constant; decisiveness fails);
(iii)~stable $+$ decisive (set $E$ to a maximally committal constant;
faithfulness fails).
\end{proposition}

\begin{proposition}[Necessity]
\label{prop:necessity}
An explanation system admits a faithful, stable, and decisive explanation
map if and only if it does \emph{not} have the Rashomon property.
\end{proposition}

\noindent
The Rashomon property is therefore the \emph{exact boundary} between
possibility and impossibility. The proofs of
Propositions~\ref{prop:tightness}--\ref{prop:necessity} are in
\emph{SI Appendix}.


%%=======================================================================
\section{Instances}
\label{sec:instances}
%%=======================================================================

We instantiate Theorem~\ref{thm:universal} in four domains chosen for
cross-disciplinary breadth. Five additional instances (attention maps,
counterfactual explanations, concept probes, model selection, language
model explanations) are detailed in \emph{SI Appendix}.

\subsection{Feature Attribution (Machine Learning)}
\label{sec:instance-attr}

Fix $P \geq 2$ features with collinearity groups. The configuration
space~$\Params$ is the set of trained models (e.g., gradient-boosted
trees, Lasso, neural networks); the observable space~$Y$ is the
input--output function; the explanation space~$\sH$ is the set of
attribution vectors $(\varphi_1, \ldots, \varphi_P)$ assigning importance
to each feature; and the incompatibility relation declares two vectors
incompatible when they rank the same collinear pair in opposite orders.
Collinearity induces a Rashomon set of near-optimal models that differ
only in which feature within a correlated group acts as the primary
predictor~\cite{damour2022underspecification, fisher2019models}. The
impossibility follows from Theorem~\ref{thm:universal}: no attribution
method (\SHAP{}, LIME, Integrated Gradients) can simultaneously be
faithful, stable, and decisive. This subsumes the result of
\citet{bilodeau2024impossibility} as a special case.

For gradient-boosted trees with per-group correlation~$\rho$, the
attribution ratio diverges as $1/(1-\rho^2)$ as $\rho \to 1$. For
Lasso at any $\rho > 0$, one correlated feature is zeroed out entirely
(infinite ratio). Empirically, SHAP rankings flip at a 33\% rate across
equivalent models on standard benchmarks.

\subsection{Causal Discovery}
\label{sec:instance-causal}

The causal discovery instance is the strongest because the Rashomon
property is \emph{derived from first principles} rather than axiomatized.

The configuration space~$\Params$ is the set of directed acyclic graphs
(DAGs) over a fixed vertex set~$V$. The observable space~$Y$ is the set
of conditional independence (CI) structures---each element records, for
every triple $(X_i, X_j \mid S)$, whether $X_i \perp\!\!\!\perp X_j
\mid S$ holds. The observation map sends each DAG to the CI structure it
encodes via $d$-separation (the Markov condition). Two DAGs map to the
same CI structure if and only if they are \emph{Markov
equivalent}~\cite{verma1991equivalence}. The explanation map is any
orientation rule that returns a fully directed graph. The incompatibility
relation declares two DAGs incompatible when they orient at least one
edge in opposite directions.

\begin{proposition}[Causal Rashomon Property---Derived]
\label{prop:causal-rashomon}
For any completed partially directed acyclic graph (CPDAG) containing at
least one undirected edge $\{X, Y\}$, there exist Markov-equivalent DAGs
$h, h'$ with opposite orientations of that edge. Hence the causal
explanation system has the Rashomon property.
\end{proposition}

\begin{proof}
An undirected edge in the CPDAG means both orientations are
Markov-consistent. Let $h$ orient it as $X \to Y$ and $h'$ as $Y \to X$.
Then $\mathsf{observe}(h) = \mathsf{observe}(h')$ (same CI structure) and
$\incomp(h, h')$ (opposite orientations).
\end{proof}

The three-variable chain $X \to Y \to Z$ and fork $X \leftarrow Y \to Z$
provide the canonical example: both encode $X \perp\!\!\!\perp Z \mid Y$
(and no other CI relation), so they are Markov equivalent, yet they
disagree on the $X$--$Y$ edge direction. No causal discovery algorithm
can distinguish them from observational data alone.

By Theorem~\ref{thm:universal}, no orientation rule can simultaneously be
faithful (return the true DAG), stable (give the same answer for
Markov-equivalent structures), and decisive (always commit to a fully
oriented graph). This is the classical Markov equivalence
underdetermination recast in our framework, with the Rashomon property
derived rather than assumed. The resolution is to report
CPDAGs~\cite{chickering2002optimal} rather than DAGs---trading
decisiveness for faithfulness and stability, exactly as the theorem predicts.

\subsection{Mechanistic Interpretability}
\label{sec:instance-mi}

Mechanistic interpretability (MI) aims to reverse-engineer neural networks
into human-understandable
circuits~\cite{elhage2022toy,conmy2023automated,anthropic2025circuit}. The
configuration space is the set of weight configurations; the observable
space is the input--output function; the explanation space is the set of
circuit decompositions; and the incompatibility relation declares two
circuits incompatible when they attribute the same computation to
contradictory subsets of network components.

\citet{meloux2025everything} found \textbf{85 distinct valid circuits}
with zero circuit error for a single XOR task, each admitting an average
of 535.8 valid interpretations. Complementary evidence from
\citet{sae_seeds_2025} shows that sparse autoencoders trained on the same
data learn different features depending on the random seed.

By Theorem~\ref{thm:universal}, no circuit explanation method can
simultaneously be faithful (reflect the actual internal circuitry),
stable (assign the same circuit to functionally equivalent networks), and
decisive (commit to a single circuit). This result does not undermine
MI---it clarifies that the goal should be to characterize the
\emph{space} of valid circuits (the ``circuit equivalence class'') rather
than to search for a single canonical decomposition. This mirrors the
CPDAG resolution for causal discovery: report which circuit features are
invariant across all valid decompositions (the ``skeleton'') and which are
decomposition-dependent (the ``reversible edges'').

For AI safety, the implication is direct: any safety argument that depends
on a specific circuit being THE mechanism of a capability (e.g.,
deception) must account for the possibility that functionally equivalent
alternative circuits exist, potentially with different safety-relevant
properties.

\subsection{Neural Degeneracy (Neuroscience)}
\label{sec:instance-neuro}

Marder and Goaillard~\cite{marder2006variability} documented that
crustacean stomatogastric ganglia producing identical motor patterns are
built from neurons with qualitatively different ionic conductance profiles.
Edelman and Gally~\cite{edelman2001degeneracy} formalized this as
\emph{degeneracy}: structurally distinct elements performing the same
function. In our framework, the configuration space is the set of
neural circuit parameterizations (ionic conductances, synaptic weights);
the observable space is the motor pattern; and the explanation space
assigns to each circuit a mechanistic account of which ion channels drive
the rhythm.

Degeneracy is precisely the Rashomon property: two circuit configurations
produce the same rhythmic output yet honestly report incompatible
mechanistic explanations (e.g., ``the rhythm is driven by the
$I_\mathrm{Ca}$ current'' versus ``the rhythm is driven by the
$I_\mathrm{h}$ current''). By Theorem~\ref{thm:universal}, no
mechanistic account of neural circuit function can simultaneously be
faithful to the specific circuit, stable across degenerate circuits, and
decisive about which ionic mechanism is primary. This connects directly to
the MI instance (Section~\ref{sec:instance-mi}): artificial and biological
neural circuits face the same structural impossibility.

\medskip

\noindent
The remaining five instances---attention maps (60\% argmax flip rate
across equivalent DistilBERT models), counterfactual explanations (23.5\%
direction flip rate), concept probes (0.90 direction instability), model
selection (80\% best-model flip rate), and language model explanations
(34.5\% citation flip rate)---are detailed with full experimental
methodology in \emph{SI Appendix}.


%%=======================================================================
\section{Resolution}
\label{sec:resolution}
%%=======================================================================

\subsection{Escape Routes}

The impossibility is tight: one desideratum must be sacrificed. Table~\ref{tab:escape-routes} summarizes the three escape routes.

\begin{table}[t]
\centering
\caption{The three escape routes from the explanation trilemma.}
\label{tab:escape-routes}
\small
\begin{tabular}{@{}llll@{}}
\toprule
Drop & Lose & Gain & Tools \\
\midrule
Faithful & May contradict & Stable + decisive & Fixed rules \\
Stable   & Changes on retrain & Faithful + decisive & \SHAP{}, single circuit \\
Decisive & Ties, uncertainty  & Faithful + stable  & \DASH{}, CPDAG \\
\bottomrule
\end{tabular}
\end{table}

\subsection{$G$-Invariant Aggregation}

Let $G$ be a group acting on the configuration space such that $g \cdot
\theta$ preserves observables: $\mathsf{observe}(g \cdot \theta) =
\mathsf{observe}(\theta)$. An explanation map is \emph{$G$-invariant} if
$E(g \cdot \theta) = E(\theta)$ for all $g, \theta$.

\begin{proposition}[$G$-Invariance Implies Stability]
\label{prop:g-stable}
If $G$ acts transitively on observe-fibers and $E$ is $G$-invariant, then
$E$ is stable.
\end{proposition}

When $G$ is finite and the explanation space admits averaging, the
canonical $G$-invariant map is the \emph{orbit average}:
\[
  \bar{E}(\theta) = \frac{1}{|G|}\sum_{g \in G} \mathsf{explain}(g \cdot \theta).
\]
This map is stable (by $G$-invariance), faithful in expectation (the
expected explanation matches the expected native explanation over the
orbit), and Pareto-optimal among stable maps. It sacrifices decisiveness:
where the orbit disagrees, the average produces ties.

The connection to classical statistics is direct. The Hunt--Stein
theorem~\cite{lehmann2005testing} establishes that among all invariant
decision procedures, the best invariant procedure is admissible. Our
$G$-invariant resolution is the explanation-theoretic instantiation: among
all stable explanation maps, the orbit average achieves the best possible
faithfulness.

\subsection{Instance-Specific Resolutions}

The abstract resolution instantiates concretely across domains:
\begin{itemize}
  \item \textbf{Feature attribution}: \DASH{}
    (\textbf{D}iversified \textbf{A}ggregation of \textbf{SH}AP)---average
    SHAP values over an ensemble of retrained models. Consensus equity
    ($\E[\varphi_j] = \E[\varphi_k]$ for symmetric features) is proved in
    Lean~4.
  \item \textbf{Causal discovery}: report the CPDAG (completed partially
    directed acyclic graph) rather than a single
    DAG~\cite{chickering2002optimal}. Undirected edges mark the
    unresolvable queries.
  \item \textbf{Mechanistic interpretability}: report the circuit
    equivalence class---the set of all valid circuits---rather than a
    single decomposition. Invariant circuit features form the ``skeleton.''
  \item \textbf{Neural degeneracy}: characterize the \emph{degenerate set}
    of ionic conductance profiles producing the target motor pattern, rather
    than committing to a single mechanistic account.
  \item \textbf{Medical diagnosis}: differential diagnosis \emph{is}
    $G$-invariant resolution---reporting the set of conditions consistent
    with the symptoms rather than committing to a single diagnosis.
\end{itemize}


%%=======================================================================
\section{Empirical Validation}
\label{sec:experiments}
%%=======================================================================

We validate the impossibility empirically across seven of the nine
instances (Table~\ref{tab:cross-instance}). For each instance, we train
multiple models that are functionally equivalent (matched on prediction
accuracy) and measure explanation instability.

\begin{table}[t]
\centering
\caption{Cross-instance empirical validation. All instability measures have
95\% bootstrap CIs excluding zero. Prediction equivalence confirms
functional identity.}
\label{tab:cross-instance}
\small
\begin{tabular}{@{}llrl@{}}
\toprule
Instance & Metric & Value & Pred.\ equiv. \\
\midrule
Attribution & SHAP flip & 33\% & $>$99\% \\
Attention & Argmax flip & 60\% & 100\% \\
Counterfactual & Dir.\ flip & 23.5\% & $\Delta$AUC$<$0.03 \\
Concept probe & $1{-}|\cos|$ & 0.90 & 97.2\% \\
Model selection & Best flip & 80\% & $\Delta$AUC$<$0.02 \\
Saliency map & Peak flip & 9.6\% & 78.8\% \\
LLM explanation & Citation flip & 34.5\% & 100\% \\
\bottomrule
\end{tabular}
\end{table}

Three results deserve emphasis. First, the 60\% attention argmax flip
rate: despite 100\% prediction agreement among 10 DistilBERT models on
SST-2, the most-attended token changes for the majority of sentences.
This is not noise---it is the Rashomon property in action. Second, the
0.90 concept direction instability: concept activation vectors extracted
from 15 equivalent MLP classifiers are nearly orthogonal (mean
$|\cos| = 0.10$), meaning the ``concept direction'' is almost entirely a
function of the random seed, not the function the model computes. Third,
the 0\% instability for negative controls (instances where the Rashomon
property does not hold, e.g., unique-solution linear regression),
confirming that the theorem's boundary condition is empirically sharp.

Full experimental details---datasets, model architectures, perturbation
protocols, confidence intervals, and negative controls---are in
\emph{SI Appendix}.


%%=======================================================================
\section{Discussion}
\label{sec:discussion}
%%=======================================================================

\subsection{Philosophy of Science: Formalizing Quine--Duhem}

The Quine--Duhem thesis~\cite{quine1951dogmas, duhem1906aim} holds that
no scientific hypothesis can be tested in isolation: empirical evidence
always underdetermines theory choice because auxiliary assumptions can
be adjusted to accommodate any observation. Van
Fraassen~\cite{vanfraassen1980scientific} elevated this to a central
argument for constructive empiricism. Despite a century of philosophical
discussion, the thesis has never been formalized as a mathematical
theorem with a precise boundary condition.

Theorem~\ref{thm:universal} provides that formalization. The configuration
space~$\Params$ corresponds to the space of theories; the observable
space~$Y$ to empirical data; the explanation map to the interpretive
framework; and the Rashomon property to the existence of empirically
equivalent theories with incompatible explanations. The theorem states:
underdetermination at the observable level propagates to underdetermination
at the explanation level, and no interpretive procedure can escape the
tradeoff. The necessity result (Proposition~\ref{prop:necessity})
sharpens the thesis: the impossibility holds \emph{if and only if} the
Rashomon property obtains, providing the exact boundary that philosophical
formulations lack.

Okasha~\cite{okasha2011theory} connected Kuhn's theory choice to Arrow's
impossibility theorem in social choice. Our result provides a
complementary connection: underdetermination of explanation is an
impossibility theorem in the same structural sense as Arrow's, with the
Rashomon property playing the role of Arrow's independence of irrelevant
alternatives.

\subsection{AI Regulation}

The EU AI Act~\cite{euaiact2024} requires ``meaningful explanations'' for
high-risk systems. SR~11-7~\cite{occ2011sr117} requires ``developmental
evidence'' documenting model decisions. Both implicitly demand
faithfulness, stability, and decisiveness simultaneously.
Theorem~\ref{thm:universal} shows this is impossible whenever models are
underspecified---which is the norm for modern machine learning systems.
The resolution is not to abandon regulation but to reframe it:
require disclosure of \emph{which} desideratum is relaxed and why,
mandate $G$-invariant explanations (e.g., ensemble-averaged attributions)
with explicit uncertainty quantification, and accept that ties where the
data are ambiguous are a feature of honest explanation, not a deficiency.

\subsection{AI Safety}

The mechanistic interpretability instance
(Section~\ref{sec:instance-mi}) has direct implications for AI safety.
If a safety argument depends on a specific circuit being the mechanism of
a dangerous capability, the impossibility theorem requires that the
argument account for alternative circuits---functionally equivalent
decompositions that may have different safety-relevant properties. The
85 valid circuits found by \citet{meloux2025everything} for a simple XOR
task suggest that real-world safety-critical capabilities will admit far
more. The constructive response is to build safety arguments on
circuit-invariant properties---features that hold across the entire
equivalence class---rather than on any single decomposition.

\subsection{Closing}

The impossibility does not counsel despair but precision: state which
property is sacrificed, aggregate optimally, and report uncertainty. The
Rashomon property is the exact boundary. When it holds, the trilemma is
inescapable. When it does not, all three properties are jointly
achievable. The task for the practitioner is to diagnose which regime
obtains and to communicate the tradeoff honestly.

%%=======================================================================
\section*{Materials and Methods}
%%=======================================================================
The Lean~4 formalization (73 files, 346 theorems, 72 axioms,
0~\texttt{sorry}) is available at [repository URL]. Experimental code and
data are in \emph{SI Appendix}. The core impossibility theorem
(\texttt{explanation\_impossibility}) depends on zero axioms beyond the
Rashomon property. Full axiom inventory and theorem list are in
\emph{SI Appendix}, Table~S1.

\section*{Acknowledgments}
[To be added.]

%%=======================================================================
\bibliography{references}
\bibliographystyle{plainnat}
%%=======================================================================

\end{document}
